\begin{theorem}[delle contrazioni o punto fisso di Banach-Caccioppoli]
\mymark{***}%
\mymargin{teor. contrazioni}%
\index{teorema!di Banach-Caccioppoli}%
\index{teorema!delle contrazioni}%
\index{punto!fisso}%
\index{contrazione}%
\index{Caccioppoli!teorema delle contrazioni}%
\index{Banach!teorema delle contrazioni}%
\label{th:banach-caccioppoli}%
Sia $X$ uno spazio metrico completo non vuoto e sia $f\colon X \to X$ una funzione $L$-lipschitziana con $L<1$ (diremo che $f$ è una \emph{contrazione}).
Allora esiste ed è unico un punto
$x\in X$ tale che $f(x) = x$.
\end{theorem}
%
\begin{proof}
\mymark{***}
Si consideri un qualunque punto $p \in X$ e si definisca
la successione $x_k\in X$ tramite la definizione ricorsiva
\[
\begin{cases}
  x_0 = p \\
  x_{k+1} = f(x_k).
\end{cases}
\]
Visto che $f$ è $L$-lipschitziana si avrà
\begin{align*}
  d(x_2, x_1) &= d(f(x_1),f(x_0)) \le L \cdot d(x_1,x_0) \\
  d(x_3, x_2) &= d(f(x_2),f(x_1)) \le L \cdot d(x_2,x_1)
  \le L^2 \cdot d(x_1, x_0) \\
  d(x_4, x_3) &= d(f(x_3),f(x_2)) \le L \cdot d(x_3,x_2)
  \le L^3 \cdot d(x_1, x_0) \\
  &\vdots
\end{align*}
possiamo quindi dimostrare induttivamente che
per ogni $m\in \NN$ si ha
\[
  d(x_{m+1}, x_m) \le L^m \cdot d(x_1, x_0).
\]
Ma allora per ogni $k\in \NN$ e per ogni $j>k$
utilizzando la disuguaglianza triangolare e facendo la somma della progressione geometrica
si ha
\[
  d(x_k,x_j) \le \sum_{m=k}^{j-1} d(x_m, x_{m+1})
   \le \sum_{m=k}^{j-1} L^m \cdot d(x_1, x_0)
   = \frac{L^k-L^j}{1-L} d(x_1,x_0).
\]
Visto che $L<1$ se $k\to +\infty$ e $j>k$ questa quantità tende a zero e quindi
risulta che $x_k$ è una successione di Cauchy. Essendo per ipotesi $X$ completo sappiamo che la successione converge $x_k \to x$ ad un punto $x\in X$.
Per la continuità di $f$, passando al limite nell'equazione
$x_{k+1} = f(x_k)$ si ottiene $x = f(x)$.
Abbiamo quindi trovato un punto fisso.
Se $y\in X$ fosse un altro punto fisso si avrebbe:
\[
  d(x,y) = d(f(x),f(y)) \le L \cdot d(x,y)
\]
che è assurdo se $L<1$ e $x\neq y$.
\end{proof}

